\section{Proofs for Section~\ref{sec:count}}\label{app:count}

This appendix proves the results of Section~\ref{sec:count}. Standard facts
from measure theory and convex analysis are cited where they are used.

\subsection{Local counts}

\begin{proof}[Proof of Lemma~\ref{lem:count:interval}]
In $F_c(v)\le F_c(v+he_i)+\delta$ the terms $c_lv_l$ with $l\ne i$ cancel,
leaving $f(v)\le f(v+he_i)+c_ih+\delta$, which is the lower bound on $c_i$.
The comparison with $v-he_i$ gives $f(v)\le f(v-he_i)-c_ih+\delta$, which is
the upper bound. The difference of the endpoints is
$\bigl(f(v+he_i)+f(v-he_i)-2f(v)\bigr)/h+2\delta/h$. If
$g(t)=f(v+te_i)-L_it^2/2$ is concave, then $g(h)+g(-h)\le2g(0)$, that is,
$f(v+he_i)+f(v-he_i)-2f(v)\le L_ih^2$.
\end{proof}

The three concentration bounds stated after
Definition~\ref{def:count:local-event} are elementary. A density bounded by
$\phi$ gives $\Prob(Y\in J)\le\phi|J|$. An interval of length $\ell$ contains
at most $\ell(M-1)/(2\sigma)+1$ points of the endpoint grid with spacing
$2\sigma/(M-1)$, so its probability under $U_{\sigma,M}$ is at most
$\ell(M-1)/(2\sigma M)+1/M\le\ell/(2\sigma)+1/M$. If
$|\Prob(Y\le t)-\nu(-\infty,t]|\le\delta_{\rm K}$ for all $t$, the same bound
holds for left limits, and the probability of every interval, which is a
difference of two such values, differs from its $\nu$-probability by at most
$2\delta_{\rm K}$.

\begin{proof}[Proof of Theorem~\ref{thm:count:local}]
Fix $v\in G$. By Lemma~\ref{lem:count:interval}, $c\in E_v^\delta$ implies
$c_i\in I_i(v,h_i,\delta)$ for every $i\in Q_v$. These intervals do not
depend on $c$ and have length at most $L_ih_i+2\delta/h_i$. By independence,
\[
 \Prob(c\in E_v^\delta)\le\prod_{i\in Q_v}\Prob\bigl(c_i\in I_i(v,h_i,\delta)\bigr)
 \le\prod_{i\in Q_v}\min\bigl\{1,\pi_i(L_ih_i+2\delta/h_i)\bigr\}.
\]
Define $\psi_i(t)=1$ for the two endpoints $t\in\{a_i,b_i\}$ and
$\psi_i(t)=\min\{1,\pi_i(L_ih_i+2\delta/h_i)\}$ for the $m_i-1$ interior
points of $G_i$. The bound reads $\Prob(c\in E_v^\delta)\le\prod_i\psi_i(v_i)$,
and summing over the tensor grid gives
$\sum_{v\in G}\prod_i\psi_i(v_i)=\prod_i\sum_{t\in G_i}\psi_i(t)$, which is the
claim. If $f$ depends on a random element independent of $c$, apply this
argument for each fixed value of that element and integrate.
\end{proof}

\begin{proof}[Proof of Corollary~\ref{cor:count:levels}]
Since $\rho_ih_{j+1}=\rho_ih_j/2$, the least admissible power of two at level
$j+1$ is $m_{ij}$ or $2m_{ij}$. Equal division into $m_{ij}$ parts is refined
by equal division into $2m_{ij}$ parts, so the partitions are nested. At
level $0$, $w_i\le\rho_iS_\rho=\rho_ih_0$, so $m_{i0}=1$. Since
$w_i/2^j\le\rho_ih_j$, $m_{ij}\le2^j$. If $m_{ij}\ge2$, minimality gives
$w_i/(m_{ij}/2)>\rho_ih_j$, that is, $h_{ij}>\rho_ih_j/2$.

Let $i$ satisfy $m_{ij}\ge2$. For every $l$, including unrefined
coordinates,
$L_lh_{lj}^2\le L_l\rho_l^2h_j^2\le c_{\rm bal}L_i\rho_i^2h_j^2
<4c_{\rm bal}L_ih_{ij}^2$. Summing over $l$,
\[
 \frac{4E_j}{h_{ij}}=\frac1{2h_{ij}}\sum_{l=1}^kL_lh_{lj}^2
 <2c_{\rm bal}kL_ih_{ij},
\]
which proves the displayed inequality.

For the counts apply Theorem~\ref{thm:count:local} with $\delta=2E_j$, so
that $L_ih_{ij}+2\delta/h_{ij}\le C_kL_ih_{ij}$ for refined coordinates;
an unrefined coordinate contributes the factor $2$. With densities,
$(m_{ij}-1)\phi_iC_kL_ih_{ij}\le\phi_iC_kL_iw_i$. With
$c_i\sim U_{\sigma_i,M}$ and $M\ge m_{ij}$,
$(m_{ij}-1)\bigl(C_kL_ih_{ij}/(2\sigma_i)+1/M\bigr)\le C_kL_iw_i/(2\sigma_i)+1$.
\end{proof}

\subsection{Corrected-corner search}

\begin{proof}[Proof of Lemma~\ref{lem:count:rounding}]
Fix $x\in C$. Let $R_1,\dots,R_k$ be independent random variables with
$R_i\in\{l_i,l_i+\eta_i\}$ and $\E R_i=x_i$; if $\eta_i=0$ put $R_i=l_i$. Let
$Y^{(0)}=x$, and let $Y^{(i)}$ be obtained from $Y^{(i-1)}$ by replacing its
$i$th coordinate with $R_i$. Every $Y^{(i)}$ lies in $C$. Conditional on
$Y^{(i-1)}$, the function $t\mapsto F_c(Y^{(i-1)}+(t-x_i)e_i)-L_it^2/2$ is
concave on $[l_i,l_i+\eta_i]$, because $F_c-f$ is linear. Jensen's
inequality gives
\[
 \E\bigl[F_c(Y^{(i)})\bigm|Y^{(i-1)}\bigr]
 \le F_c(Y^{(i-1)})+\tfrac{L_i}2\bigl(\E R_i^2-x_i^2\bigr)
 =F_c(Y^{(i-1)})+\tfrac{L_i}2\operatorname{Var}R_i
 \le F_c(Y^{(i-1)})+\tfrac{L_i\eta_i^2}8 .
\]
Iterating, $\E F_c(Y^{(k)})\le F_c(x)+\sum_iL_i\eta_i^2/8$. The point
$Y^{(k)}$ is a corner of $C$, so the least corner value is at most this
expectation.
\end{proof}

\begin{proof}[Proof of Theorem~\ref{thm:count:cells}]
Item (i) is Lemma~\ref{lem:count:rounding}, since every level-$j$ cell has
side lengths $h_{ij}$ and hence correction $E_j$.

Let $x^*$ be a global minimizer, which exists because $f$ is lower
semicontinuous and $X$ is compact. We claim that at every processed level
$j$, either $U_j=F^*$, or some processed cell $C_j\ni x^*$ of level $j$ is
not closed and is retained. At level $0$ the cell $X$ contains $x^*$. Suppose
that a processed cell $C_{j'}$ of level $j'$ contains $x^*$. If it is
closed, the closure step finds a value at most $F_c(x^*)=F^*$, so
$U_{j}=F^*$ for all $j\ge j'$, because the incumbent never increases and
never falls below $F^*$. Otherwise $\underline F(C_{j'})\le F^*\le U_{j'}$ by
(i), so $C_{j'}$ is retained and one of its children contains $x^*$ and is
processed at level $j'+1$. This proves the claim.

For (iii), if $U_j\ne F^*$, the cell $C_j$ of the claim has a corner with
value at most $F^*+E_j$ by Lemma~\ref{lem:count:rounding}, so
$U_j\le F^*+E_j$. For (ii), if $U_j\ne F^*$, the claim provides a retained
cell. For (iv), a retained cell satisfies $\underline F(C)\le U_j$, so its
least corner value is at most $U_j+E_j\le F^*+2E_j$. Such a corner $v$
satisfies $F_c(v)\le\inf_XF_c+2E_j$ and hence $c\in E_v^{2E_j}$. Assign to
every retained cell one such corner. A grid node is a corner of at most
$2^k$ cells of one level, which gives the bound on the number of retained
cells. For (v), every retained cell has $\underline F(C)\ge U_j-E_j$, because
all its corners were evaluated and $U_j$ is at most their values; hence
$U_j-\underline U_j\le E_j$. If $U_j=F^*$, then $\underline U_j\le U_j=F^*$;
otherwise the retained cell $C_j$ of the claim satisfies
$\underline F(C_j)\le F^*$. The expectation
bound follows from (iv) and Corollary~\ref{cor:count:levels}.
\end{proof}

\begin{proof}[Proof of Corollary~\ref{cor:count:approx}]
With $\rho_i=1$, $S_\rho=s$ and $h_{iJ}\le h_J=s2^{-J}$, so
$E_J\le kLs^24^{-J}/8\le\varepsilon$ by the choice of $J$. If the search stops
before level $J$, Theorem~\ref{thm:count:cells}(ii) gives an exact minimizer.
Otherwise item (v) at level $J$ gives the stated certificate with
$\underline U=\underline U_J$ and $x$ the stored incumbent point. Level $0$ uses $2^k$ evaluations. A level $j\ge1$
processes at most $2^k$ children of each cell retained at level $j-1$, and
each child has $2^k$ corners. By Theorem~\ref{thm:count:cells} the expected
number of cells retained at level $j-1$ is at most $2^k\mathcal H$, with
$\mathcal H$ from Corollary~\ref{cor:count:levels}; for the grid law,
$M\ge2^J\ge m_{ij}$ for all $j\le J$. Hence each level $j\ge1$ costs at most
$8^k\mathcal H$ expected evaluations.
\end{proof}

\subsection{Finite laws}

\begin{proof}[Proof of Lemma~\ref{lem:count:transfer}]
(a) Let $x_j=-\sigma+2\sigma j/(M-1)$. For $t\in[x_j,x_{j+1})$ with
$0\le j\le M-2$, the grid distribution function equals $(j+1)/M$, and the
continuous one lies in $[j/(M-1),(j+1)/(M-1))$. The differences are at most
$(j+1)/M-j/(M-1)=(M-1-j)/(M(M-1))\le1/M$ in one direction and less than
$(j+1)/(M-1)-(j+1)/M\le1/M$ in the other. Outside $[-\sigma,\sigma)$ both
functions are $0$ or $1$. The second claim is
Lemma~\ref{lem:count:gauss-sampler}(b), since scaling by $\sigma$ preserves the
Kolmogorov distance.

(b) For $0\le r\le n$ let
$\lambda^{(r)}=\nu_1\otimes\dots\otimes\nu_r\otimes\mu_{r+1}\otimes\dots\otimes\mu_n$.
The laws $\lambda^{(r-1)}$ and $\lambda^{(r)}$ differ only in the $r$th
factor. By Fubini's theorem,
\[
 \lambda^{(r)}(E)-\lambda^{(r-1)}(E)
 =\int\bigl[\nu_r(E_{x_{-r}})-\mu_r(E_{x_{-r}})\bigr]\,
 d\Bigl(\bigotimes_{l<r}\nu_l\otimes\bigotimes_{l>r}\mu_l\Bigr)(x_{-r}),
\]
where $E_{x_{-r}}$ is the section of $E$. By hypothesis every section is a
union of at most $C$ intervals, hence a disjoint union of at most $C$
intervals. The probability of an interval is a difference of two values of
the distribution function or of its left limits, and both differ between
$\mu_r$ and $\nu_r$ by at most $\delta_r$. The integrand is therefore at most
$2C\delta_r$ in absolute value. Summing over $r$ proves the claim.
\end{proof}

\begin{proof}[Construction and proof of Lemma~\ref{lem:count:gauss-sampler}]
Put $K=b+20$ and $e=2^{-K}$. Let $N$ be the least power of two with
$\Delta=2K/(N-1)\le e/K$, so $\log_2N=O(b)$, and let $z_j=-K+j\Delta$ for
$0\le j<N$. Write $\varphi(z)=\exp(-z^2/2)$. The weight rule below assigns
to every $j$ a number $w_j\in[0,1]$ with a common dyadic denominator $2^q$,
$q=O(b)$, and $|w_j-\varphi(z_j)|\le e/K$.

\emph{Sampler.} Repeat at most $T=\lceil4K^2\rceil$ times: draw $j$ uniformly
from $\{0,\dots,N-1\}$ using $\log_2N$ fair bits, draw $U$ uniformly from
$\{0,\dots,2^q-1\}$ using $q$ fair bits, and return $z_j$ if $U<2^qw_j$. If no
attempt succeeds, return $0$. Let $a=N^{-1}\sum_jw_j$ be the success
probability of one attempt. The output equals $z_j$ with probability
$(w_j/(Na))(1-(1-a)^T)$, so conditional on success its law $\pi$ is
proportional to the weights; the failure atom at $0$ has mass $(1-a)^T$.

\emph{Failure mass.} At least $2/\Delta-1$ grid points lie in $[-1,1]$, a
fraction of at least $(2-\Delta)/(2K+\Delta)\ge1/(2K)$ of all points, and
there $w_j\ge\varphi(1)-e/K\ge1/2$. Hence $a\ge1/(4K)$ and
$(1-a)^T\le\exp(-T/(4K))\le\exp(-K)\le e$.

\emph{Kolmogorov distance.} Let $\Phi$ be the standard normal distribution
function, and compare in four steps. First, the output law differs from
$\pi$ by at most the failure mass, hence by at most $e$ in Kolmogorov
distance. Second, let $\pi^0$ be the law proportional to
$\varphi(z_j)$. With $W=\sum_jw_j$ and $W_0=\sum_j\varphi(z_j)$, every set $A$
of grid points satisfies
\[
 \Bigl|\frac{\sum_Aw_j}{W}-\frac{\sum_A\varphi(z_j)}{W_0}\Bigr|
 \le\frac{2\sum_j|w_j-\varphi(z_j)|}{W}\le\frac{2Ne}{KW}.
\]
The Riemann estimate below gives $\Delta W_0\ge\int_{-1}^1\varphi-(K+2)\Delta\ge1$,
so $\Delta W\ge1-N\Delta e/K\ge1/2$ and the displayed quantity is at most
$4N\Delta e/K=4(2K+\Delta)e/K\le9e$. Third, $\varphi$ is bounded by $1$ and
$1$-Lipschitz. For $t\in[-K,K]$, comparing $\Delta\sum_{z_j\le t}\varphi(z_j)$
cell by cell with $\int_{-K}^t\varphi$ gives an error at most
$(2K/\Delta)\Delta^2/2+2\Delta=(K+2)\Delta$, and the same holds for the total
mass. Normalizing by masses at least $1$, the distribution function of $\pi^0$
differs from that of the normal law truncated to $[-K,K]$ by at most
$2(K+2)\Delta\le3e$. Fourth, truncation to $[-K,K]$ changes $\Phi$ by at most
$3\Phi(-K)\le3\exp(-K^2/2)\le e$. The total is at most $14e<2^{-b}$.

\emph{Weights.} Fix $j$ and put $t=z_j^2/2\le K^2/2$ and
$s=\lceil\log_2K^2\rceil+1$, so that $u=t/2^s\le1/4$. Let $P=s+3+\lceil\log_2(K/e)\rceil$.
Compute $y_0$ by rounding $\sum_{r=0}^{P}(-u)^r/r!$ to a multiple of $2^{-P}$
and clamping to $[0,1]$; the alternating series gives
$|y_0-e^{-u}|\le4^{-P-1}+2^{-P}\le2^{1-P}$. For
$r=0,\dots,s-1$, let $y_{r+1}$ be $y_r^2$ rounded to a multiple of $2^{-P}$
and clamped to $[0,1]$. Since $|\alpha^2-\beta^2|\le2|\alpha-\beta|$ on $[0,1]$
and clamping cannot increase the distance to a point of $[0,1]$, the error at
most doubles and gains $2^{-P}$ per step, so
$|y_s-\varphi(z_j)|\le2^s(2^{1-P}+2^{-P})\le e/K$. Put $w_j=y_s$ and $q=P$.
The rounded weights have $O(b)$ bits. The exact Taylor sum before rounding
has polynomial bit length: $u$ has $O(b)$ bits and $P=O(b)$, so its powers
and the factorial denominators require $O(b^2+b\log(b+1))$ bits. Thus a
weight still costs $\poly(b)$ bit operations.

\emph{Cost and support.} Each attempt uses $O(b)$ fair bits and $\poly(b)$ bit
operations, and there are $O(b^2)$ attempts. The atoms are $0$ and the
rationals $z_j=-K+2Kj/(N-1)$, of bit length $O(b)$, in $[-(b+20),b+20]$.
\end{proof}

\subsection{The growth tail}

We use three standard facts. A Lipschitz map $\R^n\to\R^n$ is differentiable
almost everywhere (Rademacher's theorem). For a Lipschitz map
$P:\R^n\to\R^n$ and a Lebesgue measurable set $A$, the area formula gives
$\int_A|\det DP(z)|\,dz=\int_{\R^n}\#(A\cap P^{-1}\{y\})\,dy$
\cite{evans2015-measure-theory-and-fine,federer1969-geometric-measure-theory}.
A finite convex function on $\R^n$ is differentiable outside a Lebesgue null
set \cite{rockafellar1970-convex-analysis}.

\begin{proof}[Proof of Theorem~\ref{thm:count:growth-tail}]
\emph{Measurability.} Let $\varepsilon>0$. If $g_*(c)\ge\varepsilon$, the
minimizer is unique and satisfies the growth inequality with the finite
coefficient $\varepsilon$, by taking limits of admissible coefficients if
necessary. For a singleton the inequality is vacuous for every finite
coefficient. Let $c_m\to c$ with $g_*(c_m)\ge\varepsilon$, and let $x_m$ be the
corresponding minimizers, so that
$f(x_m)\le f(x)+c_m^{\mathsf T}(x-x_m)-\varepsilon\norm{x-x_m}^2$ for all $x\in X$.
Along a subsequence $x_m\to\bar x\in X$, and lower semicontinuity gives
$f(\bar x)\le f(x)+c^{\mathsf T}(x-\bar x)-\varepsilon\norm{x-\bar x}^2$. Hence
$g_*(c)\ge\varepsilon$. The set $\{g_*\ge\varepsilon\}$ is closed and
$\{g_*<\varepsilon\}$ is open.

\emph{A conjugate in coefficient space.} Put $z=-c$, $\lambda=2\varepsilon$,
$\ell_i=\min_Xx_i$, $r_i=\max_Xx_i$, and
\[
 H(a)=\max_{x\in X}\bigl\{a^{\mathsf T} x-f(x)+\varepsilon\norm x^2\bigr\},
 \qquad a\in\R^n .
\]
The maximum is attained because $-f$ is upper semicontinuous and $X$ is
compact. $H$ is convex and finite. Every maximizer $x_a$ is a subgradient:
$H(a')\ge a'^{\mathsf T} x_a-f(x_a)+\varepsilon\norm{x_a}^2=H(a)+(a'-a)^{\mathsf T} x_a$.
Every subgradient $u\in\partial H(a)$ satisfies $u_i\in[\ell_i,r_i]$: for
$t>0$, $H(a+te_i)\le H(a)+tr_i$ and $H(a-te_i)\le H(a)-t\ell_i$, while
$H(a\pm te_i)\ge H(a)\pm tu_i$.

\emph{Proximal map.} Let
$P(z)=\argmin_a\{H(a)+\norm{a-z}^2/(2\lambda)\}$, which is unique by strong
convexity, and $Q(z)=z-P(z)$. Optimality gives $Q(z)/\lambda\in\partial H(P(z))$,
so $Q_i(z)\in[\lambda\ell_i,\lambda r_i]$. Monotonicity of $\partial H$ gives
$\ip{P(z)-P(z')}{Q(z)-Q(z')}\ge0$. Writing $z-z'$ as the sum of the two
differences shows that $P$ and $Q$ are monotone and $1$-Lipschitz. In
particular, for fixed $z_{-i}$, the function $t\mapsto Q_i(t,z_{-i})$ is
nondecreasing, Lipschitz, and has values in an interval of length
$\lambda w_i$.

\emph{Regular points certify growth.} Suppose $H$ is differentiable at
$a=P(z)$. Then $\partial H(a)=\{\nabla H(a)\}$, so every maximizer in the
definition of $H(a)$ equals $x^*=\nabla H(a)=Q(z)/\lambda$, and $x^*\in X$.
Maximality gives, for all $x\in X$,
$f(x)-f(x^*)-a^{\mathsf T}(x-x^*)\ge\varepsilon(\norm x^2-\norm{x^*}^2)$. Substituting
$a=z-2\varepsilon x^*$ yields
\[
 f(x)-z^{\mathsf T} x-\bigl(f(x^*)-z^{\mathsf T} x^*\bigr)\ge\varepsilon\norm{x-x^*}^2 .
\]
With $c=-z$ this is the growth inequality, so $g_*(c)\ge\varepsilon$.

\emph{The bad event has singular Jacobian.} Let $N_0$ be a Borel null set
containing the points where $H$ is not differentiable, and let
$E=P^{-1}(N_0)$, a Borel set. By the previous step,
$\{z:g_*(-z)<\varepsilon\}\subseteq E$. For every $m$, the area formula
applied to $A=E\cap[-m,m]^n$ gives
$\int_A|\det DP|=\int\#(A\cap P^{-1}\{y\})\,dy=0$, because the integrand
vanishes outside the null set $N_0$. Hence $\det DP=0$ almost everywhere on
$E$.

\emph{Trace bound.} At a point where $P$ is differentiable,
$DP+DQ=I$, and monotonicity of $P$ and $Q$ makes the symmetric parts of $DP$
and $DQ$ positive semidefinite. If $\det DP=0$, choose $v\ne0$ with $DPv=0$;
then $DQv=v$, and the trace of $DQ$, which equals the trace of its positive
semidefinite symmetric part, is at least the Rayleigh quotient
$v^{\mathsf T} DQv/\norm v^2=1$. Therefore
\[
 \ind_E(z)\le\operatorname{tr}DQ(z)=\sum_{i=1}^n\partial_iQ_i(z)
 \quad\text{for almost every }z,
\]
where outside $E$ the right side is nonnegative. Here $\partial_iQ_i$ is the
$(i,i)$ entry of $DQ$; by Fubini's theorem it coincides almost everywhere with
the derivative of the section $t\mapsto Q_i(t,z_{-i})$.

\emph{Integration.} The coordinates of $z=-c$ are independent with densities
$p_i\le\phi_i$, so null sets have probability zero. For fixed $z_{-i}$, the
section of $Q_i$ is absolutely continuous and nondecreasing with total
increase at most $\lambda w_i$, so
$\int_\R\partial_iQ_i(t,z_{-i})p_i(t)\,dt\le\phi_i\lambda w_i$. Integrating
over $z_{-i}$ and summing over $i$,
\[
 \Prob\{g_*(c)<\varepsilon\}\le\Prob(z\in E)\le\sum_i\E\,\partial_iQ_i(z)
 \le2\varepsilon\sum_i\phi_iw_i .
\]

\emph{Consequences and sharpness.} For uniform laws on $[-\sigma,\sigma]$,
$\phi_i=1/(2\sigma)$. Taking $\varepsilon=\rho/(2\sum_i\phi_iw_i)$ gives the
high-probability statement; $g_*>0$ forces a unique minimizer. For the
example, $c\ne0$ has the unique minimizer $x^*=-\sgn(c)w/2$ and
$F_c(x)-F_c(x^*)=|c|\,|x-x^*|\ge(|c|/w)|x-x^*|^2$, with equality at the other
endpoint. Hence $g_*=|c|/w$ and
$\Prob\{g_*<\varepsilon\}=\Prob\{|c|<\varepsilon w\}=\varepsilon w/\sigma$
for $\varepsilon w\le\sigma$.
\end{proof}

\subsection{Finite-law tails and the fallback budget}

We use the following block-sensitive quantifier-elimination theorem.

\begin{theorem}[Renegar {\cite[Theorem~1.1]{renegar1992-on-the-computational-complexity-and}}]\label{thm:count:renegar}
Consider a formula
\[
 Q_1x^{(1)}\cdots Q_\omega x^{(\omega)}\,
 \Psi(y,x^{(1)},\dots,x^{(\omega)})
\]
with alternating quantifier blocks of positive sizes $n_1,\dots,n_\omega$,
$\ell\ge1$ free variables $y$, and a quantifier-free matrix $\Psi$ with
$m\ge1$ atomic polynomial
sign conditions of degree at most $d\ge2$. There is an equivalent
quantifier-free formula in $y$, a disjunction of conjunctions of polynomial
sign conditions, in which the number of conjunctions, the number of atoms in
each conjunction, and their degrees are at most
$(md)^{2^{O(\omega)}\ell\prod_kn_k}$. This holds over the reals for arbitrary
real coefficients. If the coefficients are integers of bit length at most
$L$, the output coefficients have bit length at most
$(L+1)(md)^{2^{O(\omega)}\ell\prod_kn_k}$, and the formula is computed using
$(L+1)\log(L+2)\log\log(L+3)\,(md)^{2^{O(\omega)}\ell\prod_kn_k}$ bit operations and
$(md)^{O(\ell+\sum_kn_k)}$ evaluations of $\Psi$ on sign vectors.
\end{theorem}

The constant $a_{\rm R}$ of Lemma~\ref{lem:count:finite-tails} is chosen so
that, for $\omega=2$, $\ell\le2$, and $n_1=n_2=N$, the format bound is at most
$(md)^{a_{\rm R}(N+1)^2}$.

\begin{proof}[Proof of Lemma~\ref{lem:count:finite-tails}]
Let $\Phi$ be the formula describing $\mathcal D$.
If $n+m=0$, the domain is a singleton and the growth assertion is immediate;
there are no active coordinates. In the remaining argument $n+m\ge1$, so
the two quantified blocks have positive sizes.

(a) Let $\gamma(t)$ have $i$th coordinate $t$ and the fixed values elsewhere.
We claim that $g_*(\gamma(t))\ge\varepsilon$ if and only if
\[
 \exists(x,y)\ \forall(x',y'):\quad
 \Phi(x,y)\wedge\Bigl[\neg\Phi(x',y')\ \vee\
 F(x',y')+\gamma(t)^{\mathsf T} x'-F(x,y)-\gamma(t)^{\mathsf T} x-\varepsilon\norm{x'-x}^2\ge0\Bigr].
\]
If $g_*\ge\varepsilon$ with minimizer $x^*$, take $y^*$ with
$F(x^*,y^*)=f(x^*)$; then every feasible $(x',y')$ satisfies
$F(x',y')+\gamma^{\mathsf T} x'\ge f(x')+\gamma^{\mathsf T} x'\ge f(x^*)+\gamma^{\mathsf T} x^*+\varepsilon\norm{x'-x^*}^2$.
Conversely, if $(x,y)$ satisfies the formula, then taking $x'=x$ shows
$F(x,y)=f(x)$, and taking $y'$ optimal for each $x'$ gives the growth
inequality at $x$. The formula has two blocks of $n+m$ variables, one free
variable $t$, $2s+1$ atoms, and degree at most $\max\{d,2\}$, because the
terms $t(x'_i-x_i)$ and $\varepsilon\norm{x'-x}^2$ have degree two. Theorem~\ref{thm:count:renegar}
gives an equivalent quantifier-free formula in $t$ with at most $H_{\rm s}^2$
polynomial occurrences of degree at most $H_{\rm s}$. Discard the identically
zero polynomials, whose signs are constant. The others have at most
$H_{\rm s}^3$ real roots in total. Between consecutive roots and at each root
every atom has constant truth value, so the good set and its complement are
unions of at most $2H_{\rm s}^3+1$ intervals. The bound does not depend on
$\varepsilon$ or on the fixed values, which enter only as real coefficients.

(b) Under independent continuous uniform laws on $[-\sigma,\sigma]$,
Theorem~\ref{thm:count:growth-tail} gives
$\Prob\{g_*<\varepsilon\}\le S\varepsilon/\sigma$; the reduced objective $f$
is lower semicontinuous because $\mathcal D$ is compact and $F$ is
continuous. The event $\{g_*<\varepsilon\}$ is open. By (a) and
Lemma~\ref{lem:count:transfer} with $\delta_i=1/M$, the probability changes
by at most $2nC_{\rm tail}/M$. Since $\{g_*=0\}=\bigcap_{\varepsilon>0}\{g_*<\varepsilon\}$,
letting $\varepsilon\downarrow0$ gives the second bound.

(c) The same argument with the density bound $1/(\sigma\sqrt{2\pi})$ and
$\delta_i=\delta_{\rm K}$.

(d) Fix an integer assignment, a face of the continuous box (each continuous
coordinate fixed at its lower bound, fixed at its upper bound, or free), and
a coordinate $i$ fixed at a bound on that face. Let $u$ be the $k'$ free
coordinates. Condition on all coefficients other than $\gamma_i$. The free
stationarity system $\nabla_uF(u,x_{\rm fixed})+\gamma_u=0$ does not involve
$\gamma_i$, and its equations have degree at most $D$. By
Lemma~\ref{lem:sp:bezout}, which rests on the refined B\'ezout inequality
\cite{fulton1998-intersection-theory}, it has at most $D^{k'}$ nonsingular
complex solutions; for $k'=0$ there is one candidate, and if $d\le1$ and
$k'>0$ no solution is nonsingular. At each
such solution the active gradient $\partial_iF_\gamma$ equals $\gamma_i$ plus
a number that is fixed under the conditioning, so the event that its absolute
value is at most $\tau$ confines $\gamma_i$ to an interval of length $2\tau$.
Its conditional probability is at most $\tau/\sigma+1/M$ under
$U_{\sigma,M}$ and at most $\sqrt{2/\pi}\,\tau/\sigma+2\delta_{\rm K}$
under the laws of (c).

If $g_*>0$, the minimizer $x^*$ is unique. Let its integer assignment and
its smallest continuous face be the ones fixed above. Its free coordinates lie
in the open intervals, so $\nabla_uF_\gamma(x^*)=0$. For every free direction
$v$ and small $t$, $x^*\pm tv$ is feasible, and the growth inequality with a
second-order expansion gives $v^{\mathsf T}\nabla_{uu}^2F(x^*)v\ge2g_*\norm v^2$. The
free coordinates of $x^*$ therefore form one of the counted nonsingular
solutions. A union over at most $R_Z$ integer assignments, $3^{n_c}$ faces,
$n_c$ choices of $i$, and $D^{n_c}$ solutions proves the bound.
\end{proof}

\begin{proof}[Proof of Lemma~\ref{lem:count:rare-fallback}]
(a) $\E[\ind_{\mathcal E}W_1]\le\Prob(\mathcal E)\sup W_1\le B^{-1}BP(L_{\max})$.

(b) Under $U_{\sigma,M}$, Lemma~\ref{lem:count:finite-tails}(b) gives
$\Prob\{g_*<g_0\}\le\rho+2nC_{\rm tail}/M$, and part (d) gives a probability
at most $K_{\rm act}(\tau/\sigma+1/M)\le\rho+K_{\rm act}/M$ for the
active-gradient event, which contains the second part of
$\mathcal E_{\rm bad}$. The sum is at most $2\rho+(2nC_{\rm tail}+K_{\rm act})/M\le3\rho$.
Under the Gaussian-like condition the two bounds are
$\sqrt{2/\pi}\,\rho+2nC_{\rm tail}\delta_{\rm K}$ and
$\sqrt{2/\pi}\,\rho+2K_{\rm act}\delta_{\rm K}$, with sum at most $3\rho$.
\end{proof}

\subsection{The exact fallback}\label{app:count:fallback}

\begin{lemma}[One shared root for a specified tuple]\label{lem:count:shared-root}
Let $m\ge1$. For $i=1,\dots,m$, let $p_i\in\Z[t]$ have positive degree $\delta_i$, and
let a rational interval $J_i$ isolate a specified real root $a_i$ of $p_i$.
Point intervals are allowed. Let $H\ge1$ bound the coefficient and endpoint
bit lengths and put $D=\prod_i\delta_i$. A deterministic algorithm constructs
a squarefree $P\in\Z[t]$ of degree at most $D$, an interval isolating one
real root $\theta$, and rational polynomials $r_i$ of degree less than
$\deg P$ such that $a_i=r_i(\theta)$. Construction work and all coefficient
and endpoint lengths are polynomial in $D,m,H+1$, with absolute exponents.
Refinement to scalar error $2^{-q}$ costs $\poly(D,m,H+q+1)$; the same
bound gives Euclidean tuple error $2^{-q}$ after an additional
$O(\log(m+1))$ bits. Exact signs of explicitly listed fixed-degree rational
polynomials in the tuple are decidable within a polynomial in these bounds
and in the polynomial's encoding length.
\end{lemma}

\begin{proof}
\emph{A reduced tensor algebra.} Replace each $p_i$ by its primitive
squarefree part, with positive leading coefficient. Rational gcd and exact
division preserve its distinct roots and take polynomial degree/height work.
Write $d_i$ for the new degree and $N=\prod_i d_i\le D$. Divide by the
leading coefficients to work with monic rational polynomials. The algebra
\[
 \mathcal A=\Q[X_1,\dots,X_m]/(p_1(X_1),\dots,p_m(X_m))
\]
has the $N$-element basis of monomials $X_1^{e_1}\cdots X_m^{e_m}$ with
$0\le e_i<d_i$. Multiplication by $X_i$ is the tensor product of its
univariate companion matrix with the other identity matrices; denote it by
$M_i$. Over $\mathbb C$, evaluation identifies this algebra with
$\prod_{\zeta\in\mathcal T}\mathbb C$, where
$\mathcal T=\prod_i Z_{\mathbb C}(p_i)$ consists of exactly $N$ distinct
tuples. Thus the $M_i$ are simultaneously diagonalizable with joint
eigenvalues $\zeta\in\mathcal T$.

\emph{A separating form with an exact test.} For integers
\[
 0\le t\le(m-1)\binom N2,
 \qquad \lambda(t)=(1,t,\dots,t^{m-1}),
\]
use $t^0=1$, also when $t=0$. Two distinct tuples $\zeta,\eta$ have equal
projections only when the nonzero polynomial
$\sum_i(\zeta_i-\eta_i)t^{i-1}$ vanishes. Each pair excludes at most
$m-1$ integers, so the displayed list contains a form separating every
complex tuple. For each form compute
\[
 M_\lambda=\sum_i\lambda_iM_i,\qquad
 \chi_\lambda(T)=\det(TI-M_\lambda)
 =\prod_{\zeta\in\mathcal T}(T-\lambda^{\mathsf T}\zeta).
\]
The form separates precisely when
$\gcd(\chi_\lambda,\chi_\lambda')=1$. Choose the first form passing this
rational gcd test. This includes $m=1$ and $N=1$.

\emph{Coordinate maps.} For the chosen form compute independently for
each $i$
\[
 q_i(T)=-\left.\frac{\partial}{\partial u}
       \det(TI-M_\lambda-uM_i)\right|_{u=0}.
\]
The product formula gives
\[
 q_i(T)=\sum_{\zeta\in\mathcal T}\zeta_i
             \prod_{\eta\ne\zeta}(T-\lambda^{\mathsf T}\eta),
 \qquad
 q_i(\lambda^{\mathsf T}\zeta)
   =\zeta_i\chi_\lambda'(\lambda^{\mathsf T}\zeta).
\]
Since $\chi_\lambda'$ is invertible modulo $\chi_\lambda$, define
$r_i=q_i(\chi_\lambda')^{-1}\pmod{\chi_\lambda}$, taking the representative
of degree less than $N$. These identities recover every coordinate of every
Cartesian tuple. Clearing a positive denominator and removing the content
of $\chi_\lambda$ gives a squarefree integer polynomial $P$ with the same
roots. It need not be irreducible or minimal.

\emph{Selection of the designated root.} The target root is
$\theta=\sum_i\lambda_i a_i$. Every real root of $P$ corresponds to an
all-real tuple: otherwise its distinct complex conjugate would have the same
real projection, contradicting separation. Let $2^{-S}$ be an effective
strict lower bound on distances between distinct roots of $P$. Integer
polynomial separation bounds give $S\le\poly(N,H_P+1)$, where $H_P$
is the coefficient length of $P$
\cite{basu2006-algorithms-in-real-algebraic-geometry,mehlhorn2015-from-approximate-factorization-to-root}.
Put $\Lambda_0=\sum_i|\lambda_i|\ge1$. Refine each original isolator
$J_i$ to width at most $2^{-S-2}/\Lambda_0$, preserving its chosen root.
The rational interval $J_\theta=\sum_i\lambda_iJ_i$ then contains $\theta$
and has width at most $2^{-S-2}$, so it contains no other root of $P$.
Endpoint tests and closed-interval Sturm counting certify its root count.
This step selects the tuple specified by the original isolators; separation
of all Cartesian tuples alone would not select it.

\emph{Heights and work.} Companion entries have heights polynomial in
$D,H$. One common denominator is the product of the leading coefficients
of the squarefree $p_i$, of bit length polynomial in $D,m,H$. The number
of forms is $O(mN^2+1)$; tested integers $t$ have
$O(\log(m+1)+\log(N+1))$ bits, and each $\lambda_i=t^{i-1}$ has
$O(m\log((m+1)(N+1)))$ bits. Hence matrix entry lengths stay polynomial.
If $c>0$ clears their denominators, put $A_\lambda=cM_\lambda$ and
$A_i=cM_i$. The integer polynomials
\[
 \widetilde\chi(T)=\det(cTI-A_\lambda)=c^N\chi_\lambda(T),\qquad
 \widetilde q_i(T)=-[u^1]\det(cTI-A_\lambda-uA_i)=c^Nq_i(T)
\]
have coefficient lengths $O(N(H_{\rm mat}+\log(N+1)))$ by determinant
bounds. They can be computed by interpolation on $0,\dots,N$ in each of
the two variables $T,u$, using polynomially many rational determinants.
Only two-variable interpolation is involved; sample values, interpolation
denominators, and intermediate heights are polynomial in $D,m,H$.
Subresultant or Sylvester-minor bounds likewise control gcds, modular
inverses, and the $r_i$ by a polynomial with absolute exponents. The same
univariate bounds control root isolation and refinement. In particular one
may take $S=C N^2(H_P+\log(N+1)+1)$ for an effective absolute $C$, so the
source precision used to select $\theta$ is polynomially bounded.

For later evaluation, a Cauchy bound gives $|\theta|\le2^{H_P+1}$.
Termwise derivative bounds on this interval give rational
$K_i\ge\max\{1,|r_i'|\}$ of polynomial logarithmic length. Refine
$\theta$ with $q+\lceil\log_2\max_iK_i\rceil+O(1)$ bits and evaluate the
maps by rational interval arithmetic. An extra
$\lceil\tfrac12\log_2\max\{1,m\}\rceil+O(1)$ bits gives the Euclidean
bound. Finally substitute these maps in any explicitly listed fixed-degree
rational polynomial, clear positive denominators, and reduce modulo $P$.
Univariate sign determination at the selected isolated root decides its
sign, including zero, with polynomial work. This proves every assertion.
\end{proof}

\begin{proof}[Proof of Theorem~\ref{thm:count:fallback}]
\emph{A canonical minimizer.} The set $\mathcal O$ of global minimizers of
$F_\gamma$ on the compact set $\mathcal D$ is nonempty and compact. Put
$\mathcal O_0=\mathcal O$ and, for $i=1,\dots,N$, let $\mathcal O_i$ be the set
of points of $\mathcal O_{i-1}$ whose $i$th coordinate is minimal. Each
$\mathcal O_i$ is nonempty and compact, and $\mathcal O_N$ is a single point
$x^{\rm lex}$, the lexicographically least global minimizer. It exists even
when $\mathcal O$ is a continuum.

\emph{Singleton formulas.} Let $\Phi$ be $\Phi_0$ conjoined with the
disjunctions $\bigvee_{t=\ell_j}^{u_j}(x_j=t)$ for $j\in\mathcal Z$, and let
\[
 \operatorname{LexGE}(y,x)=\Bigl(\bigwedge_{j}y_j=x_j\Bigr)\vee
 \bigvee_{j=1}^N\Bigl[\Bigl(\bigwedge_{i<j}y_i=x_i\Bigr)\wedge y_j>x_j\Bigr],
\]
which has $O(N^2)$ atoms and expresses $y\ge_{\rm lex}x$. For $1\le i\le N$
consider
\[
 \Theta_i(t):\quad
 \exists x\ \forall y:\ \Phi(x)\wedge x_i=t\wedge
 \Bigl[\neg\Phi(y)\vee F_\gamma(y)>F_\gamma(x)\vee
 \bigl(F_\gamma(y)=F_\gamma(x)\wedge\operatorname{LexGE}(y,x)\bigr)\Bigr],
\]
and let $\Theta_0(t)$ be the same formula with $x_i=t$ replaced by
$F_\gamma(x)=t$. If $x$ witnesses $\Theta_i(t)$, then $x$ is feasible, no
feasible point has a smaller value, and every other minimizer is
lexicographically larger; hence $x=x^{\rm lex}$ and $t=x^{\rm lex}_i$.
Conversely, $x=x^{\rm lex}$ witnesses $\Theta_i(x_i^{\rm lex})$. Thus the truth
set of $\Theta_i$ is $\{x_i^{\rm lex}\}$, and that of $\Theta_0$ is $\{f^*\}$.
All $N+1$ formulas refer to the same point.

\emph{Elimination.} Each formula has two blocks of $N$ variables, one free
variable, at most $2s+O(N^2)$ atoms, where $s$ is the number of atoms of
$\Phi$, and degree at most $\max\{d,2\}$. Here $\log(s+1)$ is polynomial in
$I$, because the integer disjunctions contribute $\sum_j(u_j-\ell_j+1)$ atoms
with binary-encoded bounds. Multiply every atom by a positive integer that
clears its denominators. The resulting integer coefficients have bit length
$L_{\rm c}\le\poly_d(I+b)$, because each atom involves the coefficients of
$F_0$, $\Phi_0$, and at most $N$ coefficients $\gamma_i$. By
Theorem~\ref{thm:count:renegar}, an equivalent quantifier-free formula in $t$
is computed within $(L_{\rm c}+1)\log(L_{\rm c}+2)\log\log(L_{\rm c}+3)\cdot2^{\poly_d(I)}$
bit operations plus $2^{\poly_d(I)}$ evaluations of the matrix, each costing
time polynomial in its explicit length $2^{\poly_d(I)}$. The output has at most
$2^{\poly_d(I)}$ polynomials of degree at most $2^{\poly_d(I)}$ with integer
coefficients of bit length at most $(L_{\rm c}+1)2^{\poly_d(I)}$. All bounds
on the exponential factors depend only on the base data.

\emph{Univariate recovery.} Replace constant atoms by their truth values. The
singleton $\theta$ in the truth set is a root of some remaining nonconstant
polynomial: otherwise all atoms have locally constant signs near $\theta$, and
the formula would hold on a neighborhood. Let $P$ be the squarefree part of
the product of the nonconstant polynomials. Isolate the real roots of $P$.
At each root decide the sign of every atom polynomial: a polynomial vanishes
at the root exactly when the root is a root of its greatest common divisor
with $P$, which is decided by counting roots in the isolating interval; a
nonvanishing polynomial has no root in that interval, so its sign is the sign
at an endpoint. Evaluate the formula at each root; exactly one root satisfies
it. Squarefree decomposition, real-root isolation, sign determination, and
refinement of an isolating interval to width $2^{-q}$ for an integer
polynomial of degree $D$ and coefficient bit length $\tau$ take
$\poly(D,\tau,q)$ bit operations
\cite{basu2006-algorithms-in-real-algebraic-geometry,mehlhorn2015-from-approximate-factorization-to-root}.
There is therefore a base-computable $A_0=2^{\poly_d(I)}$ bounding every
scalar degree independently of $b$, while scalar coefficient and endpoint
lengths are at most $A_0(I+b+1)^{e_0}$ for a fixed exponent $e_0$.
These are compatible scalar representations of the one canonical tuple,
not yet the shared-root output.

\emph{Shared-root conversion and its base budget.} Apply
Lemma~\ref{lem:count:shared-root} to the $N+1$ scalars consisting of the
coordinates and the value. Their product degree is at most
$A_0^{N+1}$, whose logarithm is $\poly_d(I)$ and does not depend on $b$.
The lemma's absolute polynomial work and height bounds are consequently
at most $B(I+b+q+1)^{e_d}$ after enlarging the base-only integer
$B=2^{\poly_d(I)}$ and a fixed exponent $e_d$. This proves the claimed
separation of degree from coefficient height: no power of $b$ or $q$
depending on $N$ occurs. The enlarged $B$, including elimination, all
scalar recoveries, and conversion, is computed before rare-event
thresholds or sampling precision are chosen.

The source isolators select the tuple
$(x^{\rm lex},f^*)$, so at the selected root
$F_\gamma(r_1(\theta),\dots,r_N(\theta))=r_f(\theta)$.
This equality can be verified by exact sign determination there. It need
not hold modulo $P$: other Cartesian tuples may pair coordinate roots
with an unrelated value root. For an integer coordinate, refine its scalar
isolator to width less than $1/2$ and extract its unique possible integer;
the domain premise guarantees that integer exists. List it explicitly and
discard its map from the submitted representation. All other maps and the
value share $\theta$, as required by the model. Refinement with the
additional $O(\log(N+1))$ coordinate precision allowance gives (ii).

\emph{Rational feasible approximation.} Let $\mathcal D$ be a mixed box with
continuous coordinates in $[\ell_i,u_i]$, and let $G$ be the rational number
obtained by bounding every monomial derivative termwise on the box
$\prod_i[\ell_i,u_i]$; then $G\ge\max\{1,\sup\norm{\nabla F_\gamma}_1\}$ over the
box, and $\log G=\poly_d(I+b)$. Put $\delta=2^{-q}$. Refine every continuous
coordinate of $x^{\rm lex}$ to an interval of width at most
$\delta/(2G\max\{1,N\})$,
take its rational midpoint, and clip it to $[\ell_i,u_i]$; clipping does not
increase the distance to $x_i^{\rm lex}\in[\ell_i,u_i]$. Refine every integer
coordinate to an interval of width less than $1$ and take the integer in it.
The resulting rational point $y$ lies in $\mathcal D$ and has the integer
coordinates of $x^{\rm lex}$, so the segment between them lies in the box and
$0\le F_\gamma(y)-f^*\le G\norm{y-x^{\rm lex}}_\infty\le\delta/2$. Refine
$f^*$ to an interval of width at most $\delta/2$ and let $\ell$ be its lower
endpoint. Then $\ell\le f^*$ and $F_\gamma(y)-\ell\le\delta$. These refinements
also give $\norm{y-x^{\rm lex}}\le\delta$; they need
$q+O(\log(N+1))+\poly_d(I+b)$ bits and fit in the bound of (ii).
\end{proof}

\subsection{Resolution uniform over an input-length class}\label{app:count:universal}

\begin{proof}[Proof of Corollary~\ref{cor:count:universal-law}]
\emph{Grid schedules.} If $I_\Pi\le I$, monotonicity of $P_0$ gives
$M(I)\ge M_\Pi$. The schedule hypothesis supplies every-draw correctness,
the same output format and same-draw fallback, and the original structural
and numerical work factor. Only its fixed bit polynomial is now evaluated
at $I_\Pi+\log_2M(I)$, which is polynomial in $I$. The algorithm evaluates
the fixed envelope and samples by its bounded sampler; it does not
enumerate inputs, verify all promises in the class, or materialize atoms.
Certificate verification and any explicit oracle factor remain charged.

\emph{Gaussian support and precision together.} Put $H=A(I)+D(I)+21$
and $t=\lceil4\log_2H\rceil$. Since $H\ge3$,
$t\le4\log_2H+1$ and $\log_2H\le H$. With $b=A(I)+D(I)t$,
\[
 b+20\le H+4H\log_2H\le H+4H^2\le H^4\le2^t,
 \qquad 2^t\le2H^4.
\]
The integer $t$ is computed by comparing powers of two with $H^4$,
so the construction needs no real logarithm oracle.
Also $b\ge A(I_\Pi)+D(I_\Pi)t\ge b_\Pi(t)$. Thus both schedule
requirements hold for the same pair $(t,b)$: the sampler support fits the
trial box and its cap is computed for that box before drawing noise. The
integers $t,b$ are polynomially bounded in $I$, so the bit factor stays
polynomial. The schedule preserves the original structural and numerical
factor. Later $q$-bit evaluations refine the same exact output descriptor;
they neither change the law nor redraw the instance.

\emph{A supplied structural bound.} In (c) first compute the integer
envelope $F(K)$, charging its stated evaluation cost. Since the encoding
of $K$ is part of $I$, this cost is at most
$(1+F(K))^{c_F}\poly(I)$. The common resolution then dominates
each required resolution with parameter at most $K$. The sampled lengths
are bounded by $\poly(I)+(1+F(K))P_0(I)$ rather than by $\poly(I)$ alone.
Substituting this in a fixed polynomial work or output bound gives an
extra factor $(1+F(K))^c$ for a fixed $c$, enlarged to cover all charged
costs. This proves (c) while retaining its explicit dependence on $K$.
\end{proof}

We verify the schedule premises for the stated routes next. They do not
follow from fixed polynomial degree alone.

\paragraph{Effective format and height envelopes.}
Fix the degree and other format constants of a route. Explicit dimensions,
monomial lists, supplied factors, scales, curvature bounds, brackets, and
certificates are counted in its base length. Positive rationals of length
at most $I$ have logarithmic magnitudes and reciprocal magnitudes $O(I)$.
Rational matrix operations, fixed-degree derivative bounds, and rational
LP vertex bounds have polynomial bit lengths. On the domains used by the
polynomial-bit routes, this bounds the logarithms of their widths, inverse
frame bounds, curvature bounds, and reciprocal scales by effective
polynomials in $I$. These logarithmic estimates choose the law; they do
not remove large numerical ratios from expected counts.

The fallback proof separates format degree from added coefficient height.
It gives $B_\Pi\le2^{P_B(I_\Pi)}$ and work
$B_\Pi\poly_d(I_\Pi+b+q)$ with a fixed exponent, including shared-root
conversion. Integer disjunctions may have exponentially many atoms, but
their count has polynomial logarithm. Hence $\log H_{\rm s}$ and
$\log C_{\rm tail}$ in Lemma~\ref{lem:count:finite-tails}, and the
mixed-box active constant $K_{\rm act}$, have polynomial envelopes.
The fixed-block elimination formats for recourse tube constants have the
same property in the routes stated with polynomial precision. These are
effective estimates from the fixed algorithms. Finitely many estimates
can be dominated by one nondecreasing polynomial with effective integer
coefficients and degree independent of the structural parameters.

The geometric qualification matters. The compact degree-two domain
$1\le x_0\le2$, $x_i=x_{i-1}^2$ for $1\le i\le N$, has description
length $O(N\log N)$, but final width $2^{2^N}-1$. Format and coefficient
height alone do not give a polynomial logarithmic width bound there.
This example concerns the budget inference, not the difficulty of
optimizing on that particular chain. The polynomial-bit routes use boxes,
rational polytopes, fixed-depth explicit graphs, or supplied brackets
and curvature data.

\paragraph{Uniform grids and lattice termination.}
The grid closure routes select $M$ above finitely many base quantities
such as $2^J$, inverse failure budgets, fallback factors, and tail or
tube constants. The preceding envelopes bound the logarithm of each by
a fixed polynomial. In the proofs, increasing $M$ decreases the atom
term $1/M$, product-replacement errors, and finite isolation errors,
while support stays $[-\sigma,\sigma]$. Their thresholds, search box,
and cap remain valid. This verifies (a) for sparse, constrained,
continuous and native-integer recourse, and interior and bilinear
flow/TU. Aligned row noise uses the same $M(I)$ in each row, rescaled
by that row's supplied $\sigma_i$.

For Theorem~\ref{thm:qp:two}, $M\ge2n\mathsf DB$ has polynomial
logarithm, and the capped-moment proof improves with larger $M$.
Its Gaussian version similarly permits larger polynomial-bit accuracy.
For Theorem~\ref{thm:int:lowrank}, reset $J=\log_2M$: terminal curvature
error is $O(M^{-2})$, whereas value spacing is $1/[D_0(M-1)]$.
Their ratio improves with $M$ once the base inequality holds. Counts
remain valid through the new polynomial cap. Using one $M$ in every
row retains the common value denominator.

\paragraph{Gaussian-like quadratic routes.}
In Appendices~\ref{app:qp:main} and~\ref{app:qp:sep}, trial radius
$2^t\sigma$ gives nominal box width at most $2^t$ times its initial
width. The mixed-gap schedule enforces $h_J\le1$, making its gap
constant support independent. Thus $J_\Pi(t)\le J_\Pi(0)+t$; the
weaker $J_\Pi(t)\le J_\Pi(0)+2t$ would also suffice. The deterministic
node budget $(J+1)(2^J+1)^k$ and bad-event requirements give
\[
 b_{\rm req,\Pi}(t)\le A(I_\Pi)+D(I_\Pi)t
\]
for effective nondecreasing integer polynomials $A,D\ge1$. Indeed
$k\le\poly(I_\Pi)$, $J_\Pi(0)$ has a polynomial envelope, and
$\log(J+1)$ is bounded by a polynomial base term plus $t$. Rational
normalization increases encoding length only polynomially. The
two-direction requirement $2^b\ge2n\mathsf DB$ can be included by
enlarging $A$ and is independent of support. The Gaussian-weighted
counts use original projected widths, so the common trial support adds
depth and bit costs without a new numerical width factor. This verifies
(b), including anisotropic separable normalization. Original-objective
regret uses the new support $(b(I)+20)\sigma$; a calibration made with
smaller support must be recomputed.

\paragraph{Strong fields and nonlinear boundary flow.}
Theorem~\ref{thm:int:strong-field} holds at every resolution with its
actual interval probabilities $q_i(M)$ and derived $\beta(M)$. These
probabilities need not decrease under refinement: endpoint grids of
sizes $M$ and $2M$ are not nested. Hence (a) does not preserve arbitrary
distribution-dependent subcriticality. In the sufficient regime
\eqref{eq:int:sf-regime}, however, the scale premise is retained and
the lower bound on $M$ has polynomial logarithm. A common grid preserves
that sufficient regime, with $q_i,\beta$ recomputed for the new law.

For nonlinear boundary flow and its TU transfer, the available
distance-to-chart-margin proof gives
$J_\Pi,\log_2M_\Pi\le F_d(k)P_d(I_\Pi)$. This parameter function
can enter coefficient bit length, not just an analysis-only degree.
The explicit parameter powers admit a nondecreasing integer envelope in
supplied $K\ge k$ whose evaluation costs a fixed polynomial in
$1+\operatorname{size}(K)+F_d(K)$. Choose this envelope for (c).
Sampling, output, and work then carry a factor depending on $K$, with
polynomial exponents independent of $K$. Setting $K=I$ supplies a
computable input-length law but may charge that precision even when the
actual core is small. Thus the proof gives no polynomial-bit law
independent of $k$ preserving the original actual-parameter bound in
this family. This is a limitation of the available budget, not an
impossibility theorem; interior and bilinear cases belong to (a).

Finally, each oracle or curvature premise must hold on the new support.
The contracts here cover the requisite coefficient cubes or all rational
queries, rather than a single old grid. Fixing the cost model is
essential: unrelated oracle programs with arbitrary polynomial exponents
do not share a uniform work constant. These arguments standardize a
scalar family rescaled by the supplied noise scales in the theorem's
perturbation coordinates. They neither interchange law families nor
permit unbounded raw-grid refinement under a fixed atomic law.
